\documentclass[11pt]{article}
\usepackage[margin=1in]{geometry}
\usepackage{amsmath,amssymb,amsthm}
\usepackage[T1]{fontenc}
\usepackage{lmodern}
\usepackage{hyperref}
\title{There Is No 11-Arc in the Projective Plane of Order 8}
\author{quantumtool}
\date{October 19, 2026}
\begin{document}
\maketitle
\section*{Standard projective-plane interpretation}
An arc in $\mathrm{PG}(2,q)$ is a set of points with no three collinear. In a projective plane of order $q$, exactly $q+1$ lines pass through each point. If an arc has $k$ points, fix one point $P$. Each of the other $k-1$ points lies on a line through $P$, and these lines must be distinct: two other points on the same line would make three arc points collinear. Therefore $k-1\le q+1$, or $k\le q+2$.

For $q=8$, this gives $k\le 10$. There are no arcs of length 11 in $\mathrm{PG}(2,8)$, so the number of their isomorphism classes is zero, not one. The conjecture does not name an ambient projective space; this disproof uses the standard projective-plane meaning of “MDS arc,” consistent with the surrounding arc conjectures in the repository.

\section*{Formalization scope}
Lean checks the decisive cardinality contradiction: an 11-arc would require ten distinct lines through a fixed point, while an order-8 plane has only nine. The projective-plane incidence facts and the no-three-collinear injection are proved in the text and are not formalized in Lean.
\section*{Lean formalization}
Lean formalizes the final finite counting contradiction (10 distinct lines would have to fit among 9 lines). The projective-plane incidence argument supplying those counts is explicit in the report.
\end{document}
